  \begin{abstract}
  This project adds a natural-language advisory layer to an existing digital twin of a tempeh
  incubator without changing how the twin computes. The layer is a large language model, Claude
  Opus 5.5, working through wrapper tools and guided by skills, the instructions it follows. When
  the twin needs to answer a new kind of question, should the answer be coded as a new tool or
  delegated to the model over the existing tools? Inverse questions, such as how long the lid can
  stay open or when to open it, were answered by a search coded over the twin's tools and by the
  model, with skills that describe such a search, with skills that do not, and without skills.
  Answers were scored against an exhaustive search over the same tools and on whether their
  numbers came from the tools. With skills that describe the search, the model reached the
  reference in every scored run, in 53 of 54 repeats and on two questions written after the
  skills, and invented no value where none could be given. It took time of the same order as the
  coded search but 194 to 473 thousand input tokens per answer, against 130 to 138 thousand for
  one tool call. Unlike the coded search, it went to the exact limit with no margin for the
  simulation's resolution, and rarely labelled a searched value as its own in the same sentence.
  With skills that describe no search the model asked back and mostly stopped short, and without
  skills it found most limits but went past what the twin computes. For this twin, a new question
  is best delegated, with its kind described in the skill, and a search is worth coding when its
  question recurs, needs a margin for the simulation's resolution or must give the same answer
  every time.
  \end{abstract}
